\section{Time-domain simulator and its verification}\label{app:simulator}

\paragraph{Integration}
The simulator integrates Eqs.~\eqref{eq:eom}--\eqref{eq:dvf} for the state $(x,\dot x,q,\dot q)$, with the cutting force of Eq.~\eqref{eq:fcut} evaluated on the $N_t=4$ teeth and on 30 axial layers of thickness $a/30$, each represented by its mid-height. A revolution is divided into 512 steps, so that the tooth period is a whole number of steps (128), and each step is integrated by two classical fourth-order Runge--Kutta (RK4) sub-steps. The cutting force is re-evaluated at every RK4 stage: the tooth angles are advanced to the stage time, and the delayed displacement is taken from a cubic Hermite interpolant of the stored history of $x$ and $\dot x$. The usual scheme, a force computed at the start of a step and held over it, delays the force by half a step; at 512 steps per revolution it lowers the simulated stability limit by up to \SI{9}{\percent} and inflates the peak demand below \SI{1300}{\rpm} by a median factor of 1.6 (90th percentile 3.3). A run stops as divergent when $\lvert x\rvert$ exceeds \SI{1}{\centi\metre}. The entries and the kick of Section~\ref{sec:bistab-prot} are options; with all of them off the result is bit-identical to the plain step entry. The mean force of the face entry matches the integral of the slice force over the engaged arc to about $3\times10^{-3}$, and the chip of the feed ramp matches the geometric chip to $10^{-5}$.
% src: src/tdsim.c (rk4, deriv, fcut, xhist; divergence at |x| > 1e-2 m), src/tdsim.py (nsub = 2, spr = 512, nlay = 30); README §2 (held force: limit up to 9 % low; demand median 1.6x, 90th percentile 3.3x, also recomputed from results/gainmap_hold512.npz vs gainmap_main.npz below 1300 rpm: 1.64x, 3.32x), §11.2 (options, 3e-3, 1e-5); verification/c1/reference_defaults.json

\paragraph{Full surface memory}
For each of the 512 angular slots of every layer the simulator stores the step of the deepest pass; a pass that cuts is deeper than every earlier pass at that slot and replaces it. The chip of Eq.~\eqref{eq:chip} becomes the smaller of the chip against the last pass ($k=1$) and the chip against the deepest earlier pass at the two slots that bracket the stage angle ($k$ tooth periods back),
\begin{equation}\label{eq:surf}
  h_j=\min_{k}\Bigl[k\,f_z\sin\phi_j-\bigl(x(t)-x(t-kT)\bigr)\cos\phi_j\Bigr],
\end{equation}
and material that no earlier tooth has cut, as during the face entry, gives the fresh chip $f_z\sin\phi_j-x(t)\cos\phi_j$. The scheme was compared with a brute force that takes the least chip over the last 600 passes, which is the full memory for runs of up to 600 passes. A stable cut agrees to \SI{0.32}{\percent}, whereas single regeneration is \SI{4.9}{\percent} off, because the tool leaves the cut briefly during the entry transient. Clipped chatter is chaotic and separates pointwise, but its amplitude agrees to \SI{0.4}{\percent}, and the same run diverges (\SI{1050}{\rpm}) in both. With only the last four passes remembered, two of four chatter cases diverge where the full memory stays bounded, and a third is \SI{27}{\percent} too large.
% src: src/tdsim.c (fcut_multi, deep[], full); verification/tdsim/check9_full_memory.json (stable: 0.0032 vs brute, 0.0486 vs single; chatter amplitudes 0.5896/0.5917 and 0.4345/0.4352 mm; nreg 4: 2 diverge, 0.554 vs 0.435 mm)

\paragraph{Convergence}
The reference is the exact zero-order-hold solution of the linear loop at 8192 and 16384 steps per revolution, extrapolated to zero step (Richardson) and accurate to about \SI{0.04}{\percent} (five cases, ten revolutions each). At 512 steps per revolution the scheme deviates from it by \SI{0.74}{\percent} at \SI{500}{\rpm} ($a=\SI{1.5}{\milli\metre}$, $g=0$), by \SI{0.07}{\percent} at \SI{1950}{\rpm}, and by no more than the accuracy of the reference in the three other cases; the held-force scheme deviates by \SIrange{1.8}{8.3}{\percent}, and the 8192-step solution without extrapolation by \SIrange{0.1}{0.6}{\percent}. The peak demand at $a_{\lim}$ agrees between 512 and 8192 steps to \SI{0.18}{\percent}.
% src: verification/tdsim/check7_final_scheme.json (trajectory: rel_err_fmode1 0.0074, 0.00036, 0.000034, 0.00011, 0.00065; rel_err_fmode0 0.0176-0.0831; rel_err_zoh8192 0.0010-0.0057; rel_err_fmode1_16384 <= 0.00039; demand |rel| <= 0.0018)

\paragraph{Stability limit and layers}
An independent first-order semi-discretisation \citep{insperger2002semi,insperger2004updated,insperger2011book} with Richardson extrapolation, written separately from the code of Section~\ref{sec:model}, was compared with the ZOA limit at 12 speeds, four gains (0, 300, 800 and \SI{1500}{\volt\second\per\metre}) and both damping ratios (96 cases). With the coefficient averaged over the tooth period, which is the model of the ZOA, the largest relative difference is $2\times10^{-5}$; against a classical lobe sweep on a \SI{1}{\rpm} grid it is $6\times10^{-6}$, with no missed roots; with the time-periodic coefficient $aH(t)$ of Eq.~\eqref{eq:aHt} for the \SI{45}{\degree} helix, the ZOA is conservative by \SIrange{0}{1.04}{\percent}. With 30 layers the settled demand exceeds the closed form of Eq.~\eqref{eq:xp} by up to \SI{8.7}{\percent} at low speeds (Section~\ref{sec:feas-new}), because the layers are thicker than the axial wavelengths of the higher force harmonics; with 240 layers the difference falls within \SI{0.8}{\percent} at the three speeds checked. The layers change the peak entry demand of Eq.~\eqref{eq:pub} by $-0.03$ to $+\SI{0.85}{\percent}$ (four speeds), and their effect on the stability limit is bounded by the \SIrange{0.08}{0.73}{\percent} by which the long-run onsets below exceed the ZOA limit.
% src: verification/zoa_sdm/README.md §3-4, verification/zoa_sdm/current_grid_check.json (2.0e-5, 6.0e-6), verification/zoa_sdm/compare_12speeds.csv; README §3 (0-1.05 %); verification/feasibility/cf_vs_td_layers.json (settled, entry_peak); verification/tdsim/check8_onset_long_runs.json

\paragraph{Chatter classifier}
A settled cut converges to a $T$-periodic orbit. The classifier D1 samples $x$ once per tooth period and examines the change $\Delta x_k$ between successive samples over the second half of the window, divided into 25 blocks. A run is stable when the per-tooth growth factor fitted to the logarithm of the block RMS of $\Delta x_k$ is below $1-2\times10^{-5}$, the largest RMS of the last five blocks is below 0.98 times that of the first five, the final RMS is below that of the entry transient, and $\lvert x(t)-x(t-T)\rvert$ stays below $0.25f_z$ over the last \SI{20}{\percent} of the window, or when $\Delta x_k$ decays to round-off; a limit cycle of any size counts as chatter. After a step entry, a 300-revolution window places the time-domain limit \SIrange{0.36}{2.2}{\percent} above the ZOA limit in eight cases, whereas runs of 2000 revolutions place the onset only \SIrange{0.08}{0.73}{\percent} above it. Just above the onset, chatter saturates on a small limit cycle that the cut needs more than 300 revolutions to reach, so the shorter window biases these boundaries high by \SIrange{0.05}{1.5}{\percent}. Near the forced resonance (\SIrange{1970}{2030}{\rpm}, small gains) the 300-revolution boundaries reach 1.015--$1.05a_{\lim}$, whereas the onset lies at 1.004--$1.012a_{\lim}$, a bias of \SIrange{0.8}{4.4}{\percent}. Kick-based limits and clipped limits well below $a_{\lim}$ are not affected.
% src: src/tdboundary.py (classify_run: delta 2e-5, 25 blocks, q < 0.98, r_end < r_first, dx_end < 0.25 fz); verification/tdsim/check7_final_scheme.json (boundary: +0.36 % to +2.20 %), check8_onset_long_runs.json (onset +0.08 % to +0.73 %; bias 0.05-1.47 %), check10_forced_resonance_long_runs.json (td300 1.0149-1.05, onset 1.0039-1.0121)
